%%%%%%%%%%%%%%%%%%%%%%%%%%%%%%%%%%%%%%%%%%%%%%%%%%%%%

\begin{frame}
    \frametitle{Metabolische Flussanalyse (MFA)?}

\textbf{Ziel:} Bestimmung der Reaktionsraten (Flüsse) in einem Stoffwechselnetzwerk

\textbf{Grundidee:}
\begin{itemize}
  \item Jedes Stoffwechselnetzwerk besteht aus Metaboliten und Reaktionen
  \item Stationärer Zustand: Konzentration jedes Metaboliten bleibt konstant
\end{itemize}

\textbf{Mathematische Formulierung:}
\begin{itemize}
  \item $S \in \mathbb{R}^{m \times n}$: Stöchiometriematrix (m Metabolite, n Reaktionen)
  \item $v \in \mathbb{R}^{n}$: Vektor der Flussraten
  \item \textbf{Nebenbedingung:} $S \cdot v = 0$ (Konzentration bleibt konstant)
  \item \textbf{Zielfunktion:} z.\,B. $c^T v \rightarrow \max$ (z.\,B. Wachstum)
\end{itemize}
\end{frame}



%%%%%%%%%%%%%%%%%%%%%%%%%%%%%%%%%%%%%%%%%%%%%%%%%%%%%
\begin{frame}
    
    \frametitle{Beispiel: Vereinfachter zentraler Kohlenstoffmetabolismus}

\textbf{Reaktionen:}
\begin{itemize}
  \item $v_1$: Glukose\textsubscript{ext} $\rightarrow$ Glukose-6-P
  \item $v_2$: Glukose-6-P $\rightarrow$ Pyruvat
  \item $v_3$: Pyruvat $\rightarrow$ 2 ATP
  \item $v_4$: Glukose-6-P $\rightarrow$ 1 ATP (alternativer Weg)
  \item $v_5$: ATP $\rightarrow$ Biomasse (Zielreaktion)
\end{itemize}

\vspace{0.5em}
\textbf{Metabolite:} Glukose-6-P, Pyruvat, ATP

\vspace{0.5em}
\textbf{Stöchiometriematrix $S$:}

\[
\begin{array}{c|ccccc}
 & v_1 & v_2 & v_3 & v_4 & v_5 \\
\hline
\text{Glu-6-P} & +1 & -1 &  0 & -1 & 0 \\
\text{Pyruvat} &  0 & +1 & -1 &  0 & 0 \\
\text{ATP}     &  0 &  0 & +2 & +1 & -1 \\
\end{array}
\]
\end{frame}

% Folie 2: Mathematische Formulierung
\begin{frame}{Mathematische Formulierung}

\textbf{Nebenbedingungen (Steady-State):}
\[
\begin{aligned}
v_1 - v_2 - v_4 &= 0 \quad \text{(Glu-6-P)} \\
v_2 - v_3 &= 0 \quad \text{(Pyruvat)} \\
2v_3 + v_4 - v_5 &= 0 \quad \text{(ATP)}
\end{aligned}
\]

\textbf{Zielfunktion:}
\[
\max_{v} \; c^T v = \max v_5
\]

\textbf{Interpretation:} Maximierung der ATP-Nutzung für Wachstum unter Massenerhaltung

\textbf{Frage:} Wie kann das LGS gelöst werden, unter Einhaltung der Bedingung? 
$\Rightarrow$ {\bf Lagrange-Multiplikatoren}

\end{frame}

